\section{The one-decision problem and the online checkpoint rule}
\label{sec:b-problem}

\paragraph{Population.}
The population is a distribution over \emph{templates} --- task specifications
that generate one or more \emph{roots} (independent runs) and, along each root,
one eligible \emph{checkpoint}.  The template is the independent unit; roots
generated from one template share its structure and are never split across
development, calibration and test.  A root that never reaches an eligible
checkpoint contributes a reachability record and no branch; it is reported, not
silently dropped.  Conditional checkpoint-population value is not automatically
whole-task value, and the reachability rate is a reported quantity.

\paragraph{The decision.}
At the checkpoint the controller observes pre-treatment information
$W=(Z,M)$: a registered representation $Z=f(H)$ of the recorded history and a
pre-treatment feasibility mask $M=(A_g)_{g\in\Gset}$ over the declared action
set $\Gset=\{g_0,g_v\}$ (an optional recovery action $g_r$ is discussed in
\cref{sec:b-bridge} and is not part of the first study).  It chooses one action
and then the frozen continuation $\rho$ runs to the registered horizon.  The
decision is made once per root; learning a stopping time, or choosing when to
intervene at every checkpoint, is outside this paper's scope.

\paragraph{Online-reachable checkpoint rule.}
The checkpoint is selected by a predeclared rule that is executable from
information available at the time: \emph{the first completed candidate
evaluation on the root that satisfies fixed structural criteria}, evaluated
before any intervention.  The rule must not depend on the eventual trajectory's
future checkpoint window; a rule that selects uniformly from a window defined
by later steps cannot be implemented online and is not used.  If a fixed
externally sampled step is preferred, it is sampled \emph{before} the trajectory
and the handling of roots that never reach it is declared in advance.  The
rule's version is recorded on every branch record, and a change to it is a
protocol revision.

\paragraph{Conditions and mixing.}
Roots are assigned to one of the registered conditions --- a verifiable and
repairable configuration mismatch; a fault that is not actionable under the
particular verification response; a latent semantic inconsistency despite
successful execution; and no injected fault --- by a mechanism whose assignment
is independent of the monitor score.  The condition label is
\emph{evaluator-only}: it never enters the feature construction, the learner's
inputs, or the controller's information.  No-injected-fault is not guaranteed
error-free; actual agent behaviour under it is diagnosed separately.  Artificial
mechanisms establish controlled differences and not prevalence in natural
deployments; unmodified tasks are included where practical, and if only
controlled faults support the result, the paper is framed accordingly.

\paragraph{Why the problem is not a stopping problem.}
The one-step restriction is what makes the decision estimand well posed: the
counterfactual continuation after each action is generated by the same frozen
$\rho$, so a branch contrast identifies the action's effect at the checkpoint
rather than a mixture of the action and a policy that would have adapted after
it.  Sequential choice over checkpoints is a different problem with a different
positivity requirement; \cref{sec:b-bridge} states what it would add and why it
is not claimed here.
